\section{Experiments}
\label{sec:experiments}


\begin{figure}[t]
    \centering
    \includegraphics[width=\linewidth]{Images/simulation_tasks.png}
    \caption{The four RoboTwin tasks: Beat Hammer, Hanging Mug, Open Microwave, and Put Object into Cabinet.}
    \label{fig:simulation_tasks}
\end{figure}

\begin{figure}[t]
    \centering
    \includegraphics[width=\linewidth]{Images/real_world_experiment_setup.png}
    \caption{Real-world setup on the DOBOT X-Trainer dual-arm platform. The bottom row distinguishes demonstration objects (Training Set) from held-out evaluation objects (Test Set), which differ in geometry and appearance.}
    \label{fig:real_world_experiment_setup}
\end{figure}

We conduct a series of experiments to evaluate the proposed representation from both representation and manipulation perspectives. We aim to answer the following questions:
\begin{enumerate}
    \renewcommand{\labelenumi}{(\arabic{enumi})}
    \item \textbf{Representation Robustness:} Do feature visualizations exhibit part discrimination and consistent patterns across object instances in simulated and real observations?
    \item \textbf{Manipulation Effectiveness:} Does the resulting representation improve policy performance across the evaluated simulation tasks and held-out real-world object instances?
    \item \textbf{Component Contributions:} How do support/query separation and the three field-learning objectives affect downstream manipulation success?
\end{enumerate}

\subsection{Experimental Setup}

\paragraph{Simulation setup.}
We evaluate four RoboTwin tasks~\cite{chen2025robotwin} spanning single-arm and dual-arm manipulation: \emph{Hanging Mug} hangs a mug on a rack through its handle; \emph{Beat Hammer} grasps the handle and strikes a target with the hammer head; \emph{Open Microwave} pulls the door handle; and \emph{Put Object into Cabinet} grasps the cabinet handle and opens the door. These tasks require localization of functional parts rather than object identity alone (Fig.~\ref{fig:simulation_tasks}). We follow the standard RoboTwin data split without modifying the benchmark data partitioning or processing protocol. All methods use the same demonstrations and test instances. Each policy is trained for 3,000 epochs on 50 demonstrations and evaluated for 100 episodes per seed.


\paragraph{Real-world setup.}
We further evaluate four bimanual tasks on a DOBOT X-Trainer dual-arm platform equipped with two ZED 2i RGB-D cameras, as shown in Fig.~\ref{fig:real_world_experiment_setup}. For each task, we collect 50 teleoperated demonstrations using designated training objects and evaluate on held-out instances with different geometry and appearance. The tasks are (1) \emph{Grasp Mug}, grasping the mug handle and placing the mug stably on the target base; (2) \emph{Beat Cube}, grasping the hammer handle and striking the cube with the hammer head; (3) \emph{Stir Mug}, grasping the mug handle with the left hand and the spoon handle with the right hand, then stirring inside the mug with the spoon head; and (4) \emph{Pour Water}, grasping the kettle handle and pouring water from its spout into the target mug. These functional-part requirements define task success. Each method is evaluated for 20 trials per task under matched initial conditions, and we report the success rate.


\paragraph{Baselines.}
\textbf{DP3}~\cite{ze20243d} is our primary geometry-only baseline. \textbf{DINOv2 Lifting} is a controlled baseline inspired by G3Flow~\cite{chen2025g3flow} and D$^3$Fields~\cite{wang2023d}: frozen DINOv2~\cite{oquab2023dinov2} descriptors are sampled at image projections of observed object points and lifted into 3D. \textbf{Utonia Point-wise} attaches frozen Utonia~\cite{zhang2026utonia} descriptors directly to observed points, without field construction or resampled queries. We additionally compare with \textbf{G3Flow}~\cite{chen2025g3flow} in simulation, using its complete semantic-flow pipeline rather than treating DINOv2 Lifting as its reproduction.


\paragraph{Implementation and evaluation details.}
Grounded SAM~\cite{ren2024groundedsam} segments task-relevant objects. Masked depth pixels are back-projected from calibrated views, transformed into the robot/world frame, and fused into object clouds. DINOv2 Lifting, Utonia Point-wise, and ours share these masks and reconstruction steps, as well as a DP3-based interface: separate PointNets encode the scene and semantic object clouds, and their features are concatenated with proprioception. The policy objective and action representation are unchanged. Point budgets follow Sec.~\ref{sec:method_policy}.

For G3Flow, we follow the official implementation and deployment procedure, including object-model-based semantic-flow construction and FoundationPose tracking to update feature locations as objects move. Unlike per-observation DINOv2 Lifting, this pipeline maintains object-aligned features through pose tracking. We retain these method-specific components rather than forcing G3Flow into the controlled lifting interface.

All comparisons use matched task demonstrations and evaluation episodes. DP3, DINOv2 Lifting, Utonia Point-wise, and ours additionally share observation histories and training schedules. Policies are trained for 3,000 epochs with batch size 128 on a single NVIDIA RTX 4090 GPU, also used for inference. Simulation reports mean success and standard deviation over training seeds, with 100 episodes per seed; real-world results report successes out of 20 trials per task. 





\subsection{Simulation Experiments}
\label{sec:simulation_results}

\begin{table}[t]
\centering
\scriptsize
\setlength{\tabcolsep}{1.0pt}
\captionsetup{justification=centering,labelsep=newline,font=footnotesize}
\caption{\textsc{Comparative Evaluation on Four RoboTwin Tasks (Success Rates in \%).}}
\label{tab:simulation_results}
\begin{tabular*}{\columnwidth}{@{\extracolsep{\fill}}l|cccc|c@{}}
\toprule
\textbf{Method}
& \makecell{\textbf{\textit{Hang}}\\\textbf{\textit{Mug}}}
& \makecell{\textbf{\textit{Beat}}\\\textbf{\textit{Hammer}}}
& \makecell{\textbf{\textit{Open}}\\\textbf{\textit{Micro.}}}
& \makecell{\textbf{\textit{Put}}\\\textbf{\textit{Cabinet}}}
& \textbf{\textit{Average}} \\
\midrule
DP3~\cite{ze20243d} & 23.0\stdval{4.0} & 72.0\stdval{2.0
} & 62.0\stdval{3.0} & 72.0\stdval{2.0} & 57.3\\
DINOv2 Lifting & 23.0\stdval{7.0} & 78.0\stdval{5.0} & 58.0\stdval{3.0} & 68.0\stdval{4.0} & 56.8 \\
Utonia Point-wise & 31.0\stdval{4.0} & 77.0\stdval{2.0} & 65.0\stdval{1.0} & 76.0\stdval{1.0} & 62.3 \\
G3Flow~\cite{chen2025g3flow} & 25.0\stdval{5.0} & 78.0\stdval{5.0} & 64.0\stdval{2.0} & 78.0\stdval{3.0} & 61.3 \\
\midrule
\textbf{Ours} & \textbf{37.0}\stdval{3.0} & \textbf{84.0}\stdval{1.0} & \textbf{73.0}\stdval{1.0} & \textbf{83.0}\stdval{1.0} & \textbf{69.3} \\
\bottomrule
\end{tabular*}
\end{table}



Table~\ref{tab:simulation_results} reports the controlled representation comparisons and the complete G3Flow baseline. Our method achieves the highest success rate on every task, averaging 69.3\%: 12.0 percentage points above DP3, 7.0 above Utonia Point-wise, and 8.0 above G3Flow. These gains across handle localization, tool-part discrimination, and articulated interaction support the utility of the learned semantic cues.

DINOv2 Lifting improves Beat Hammer but matches or underperforms DP3 on the other three tasks, yielding an average of 56.8\%. Its larger standard deviations on several tasks indicate greater variability across training runs; these statistics alone do not identify multi-view fusion as the cause. Utonia Point-wise improves the DP3 average by 5.0 percentage points, while our learned field adds a further 7.0 points. These comparisons support the utility of the learned semantic representation, but do not separately identify the effects of field construction and additional part supervision. The ablations below examine the learning objectives and query sampling design.



\paragraph{Comparison with G3Flow.}
G3Flow averages 61.3\%, improving on DINOv2 Lifting by 4.5 percentage points. It improves three tasks and matches Beat Hammer, but remains below ours on all four. Its standard deviations (2--5 points) are higher than ours (1--3 points) on every task, although no higher than DINOv2 Lifting. During evaluation, we observed drift in FoundationPose estimates, which can displace object-aligned semantic features and destabilize policy inputs. This is a plausible contributor to the remaining variability, not an isolated causal finding. Our field instead reads semantics from each current support observation without propagating tracked object poses. The comparison therefore evaluates complete pipelines rather than isolating the effect of tracking alone.

\subsection{Real-World Experiments}
\label{sec:real_results}

\begin{figure}[t]
    \centering
    \includegraphics[width=\linewidth]{Images/real_world_experiment.png}
    \caption{Real-world manipulation examples on held-out instances: Beat Cube, Grasp Mug, Pour Water, and Stir Mug.}
    \label{fig:real_world_tasks}
\end{figure}



\begin{table}[t]
\centering
\scriptsize
\setlength{\tabcolsep}{1.0pt}
\captionsetup{justification=centering,labelsep=newline,font=footnotesize}
\caption{\textsc{Comparative Evaluation on Four Real-World Tasks (Success Rates in \%).} Each method is evaluated for 20 trials per task.}
\label{tab:real_results}
\begin{tabular*}{\columnwidth}{@{\extracolsep{\fill}}l|cccc|c@{}}
\toprule
\textbf{Method}
& \makecell{\textbf{\textit{Grasp}}\\\textbf{\textit{Mug}}}
& \makecell{\textbf{\textit{Beat}}\\\textbf{\textit{Cube}}}
& \makecell{\textbf{\textit{Stir}}\\\textbf{\textit{Mug}}}
& \makecell{\textbf{\textit{Pour}}\\\textbf{\textit{Water}}}
& \textbf{\textit{Average}} \\
\midrule
DP3~\cite{ze20243d} & 35.0 & 40.0 & 15.0 & 20.0 & 27.5 \\
DINOv2 Lifting & 35.0 & 45.0 & 25.0 & 25.0 & 32.5 \\
Utonia Point-wise & 35.0 & 45.0 & 35.0 & 25.0 & 35.0 \\
\midrule
\textbf{Ours} & \textbf{85.0} & \textbf{85.0} & \textbf{50.0} & \textbf{50.0} & \textbf{67.5} \\
\bottomrule
\end{tabular*}
\end{table}

Figure~\ref{fig:real_world_tasks} shows representative task executions on held-out objects. Table~\ref{tab:real_results} reports success under the matched evaluation protocol. Our method achieves the highest success rate on every task and an average of 67.5\%, compared with 35.0\% for Utonia Point-wise, 32.5\% for DINOv2 Lifting, and 27.5\% for DP3. This corresponds to gains of 32.5 and 40.0 percentage points over the strongest semantic baseline and the geometry-only baseline, respectively. The advantage is most pronounced on \emph{Grasp Mug} and \emph{Beat Cube}, where our method succeeds in 17 of 20 trials, compared with 7 and 9 trials for Utonia Point-wise.

The margin over Utonia Point-wise is larger in the real world than in simulation (32.5 versus 7.0 percentage points). Depth noise, partial visibility, calibration residuals, and segmentation errors may contribute by destabilizing point-attached features. Our field is designed to reduce such variation through object-conditioned readout. However, the evaluations use different tasks and object distributions, so the gap cannot be attributed solely to sensing noise. Section~\ref{sec:representation_stability} examines relevant qualitative feature patterns.


\subsection{Representation Analysis}
\label{sec:representation_stability}

We qualitatively examine whether the representations preserve both \emph{intra-instance part discrimination} and \emph{cross-instance part consistency}. Feature embeddings are projected to RGB with PCA for visualization. For each method and object category, a single PCA basis is fitted jointly to all displayed instances and then applied without instance-specific normalization. Colors are therefore comparable across instances within the same method-category group: consistent colors at corresponding functional parts indicate aligned embeddings, whereas color changes across instances reveal representation drift. Absolute colors should not be compared across different methods or categories because their PCA bases are fitted independently.


\paragraph{Simulation representation.}


\begin{figure}[t]
    \centering
    \includegraphics[width=0.8\linewidth]{Images/simulation_feature.png}
    \caption{Part-aware field embeddings queried on simulation instances from five object categories. Within each category, all instances share one PCA projection. Corresponding functional parts exhibit consistent colors across instances, while distinct parts remain clearly separable.}
    \label{fig:simulation_feature}
\end{figure}


Figure~\ref{fig:simulation_feature} visualizes embeddings across five categories with substantial shape variation. Corresponding regions retain similar colors: mug handles remain distinct from bodies, and hammer heads from handles. The other categories exhibit analogous part distinctions. These patterns suggest cross-instance consistency while preserving functional detail within each object.

\begin{figure}[t]
    \centering
    \includegraphics[width=\linewidth]{Images/real_world_feature.png}
    \caption{Cross-instance feature comparison on real observations of mugs and hammers. DINOv2 lifting is view-dependent, and point-wise 3D descriptors remain sensitive to observed sampling. Our field preserves aligned part-level patterns across geometrically and visually different instances.}
    \label{fig:real_world_feature}
\end{figure}

\paragraph{Real-world representation.}
Figure~\ref{fig:real_world_feature} compares DINOv2 Lifting, Utonia Point-wise, and our field under physical RGB-D observations. DINOv2 produces meaningful local structure within an individual observation: it can often separate a mug handle from its body or a hammer head from its handle. However, the same functional part changes color markedly across instances and viewpoints, and some regions exhibit abrupt, spatially nonuniform responses. Thus, its descriptors are locally discriminative but not reliably aligned across objects of the same category after multi-view lifting. Utonia Point-wise reduces the view dependence of the 2D features and preserves recognizable geometry, yet corresponding parts still exhibit color variation across the displayed instances. These visualizations do not isolate whether that variation is caused by shape, sampling, or other observation differences.

Our field retains coherent handle/body and head/handle patterns across the four mug and four hammer instances despite geometry and appearance differences. This combines part discrimination with cross-instance consistency, in agreement with the policy gains in Table~\ref{tab:real_results}. Controlled perturbation tests are still needed to establish the contribution of sensing noise.



\subsection{Ablation Study}
\label{sec:ablation}

We isolate the contribution of the support/query separation and the three objectives used for semantic-field learning. All variants retain the frozen Utonia encoder, DP3 policy backbone, demonstrations, point budget, and optimization settings. The coupled support/query variant preserves the complete field and training objective, but selects its 256 policy-facing query coordinates directly from the sampled support set instead of independently resampling them from the observed object cloud. The remaining variants retain the full support/query-decoupled architecture and each remove exactly one loss from Eq.~\eqref{eq:total_loss}. We report the average success rate across the four simulation tasks.

\begin{table}[t]
\centering
\scriptsize
\setlength{\tabcolsep}{1.0pt}
\captionsetup{justification=centering,labelsep=newline,font=footnotesize}
\caption{\textsc{Ablation Study of Representation Design and Field-Learning Objectives.} Average success rate (\%) across four simulation tasks.}
\label{tab:ablation}
\begin{tabular*}{\columnwidth}{@{\extracolsep{\fill}}l|c|ccc|c@{}}
\toprule
\textbf{Variant}
& \makecell{\textbf{S/Q}\\\textbf{Sep.}}
& $\boldsymbol{\mathcal{L}_{\mathrm{ce}}}$
& $\boldsymbol{\mathcal{L}_{\mathrm{con}}}$
& $\boldsymbol{\mathcal{L}_{\mathrm{cr}}}$
& \makecell{\textbf{\textit{Average}}\\\textbf{\textit{(\%)}}} \\
\midrule
Coupled support/query &  & $\checkmark$ & $\checkmark$ & $\checkmark$ & 66.7 \\
Without $\mathcal{L}_{\mathrm{ce}}$ & $\checkmark$ &  & $\checkmark$ & $\checkmark$ & 61.3 \\
Without $\mathcal{L}_{\mathrm{con}}$ & $\checkmark$ & $\checkmark$ &  & $\checkmark$ & 64.7 \\
Without $\mathcal{L}_{\mathrm{cr}}$ & $\checkmark$ & $\checkmark$ & $\checkmark$ &  & 68.9 \\
\midrule
\textbf{Full model} & $\checkmark$ & $\checkmark$ & $\checkmark$ & $\checkmark$ & \textbf{69.3} \\
\bottomrule
\end{tabular*}
\end{table}

Table~\ref{tab:ablation} reports 69.3\% for the full model and 66.7\% for coupled support/query, a 2.6-point difference. Both variants retain field-based decoding and all three objectives, so this comparison tests query sampling, not a field against a point-wise model with matched part supervision.

Removing $\mathcal{L}_{\mathrm{ce}}$ reduces success to 61.3\%, the largest drop (8.0 points), supporting explicit functional-part anchoring. Removing $\mathcal{L}_{\mathrm{con}}$ yields 64.7\%, a 4.6-point drop, supporting alignment of corresponding parts across instances. Without $\mathcal{L}_{\mathrm{cr}}$, success is 68.9\%; this 0.4-point difference does not establish a reliable benefit without variability estimates. These ablations measure downstream success, not cross-view feature consistency. Controlled perturbation tests are needed to assess robustness directly. Part anchoring and alignment thus provide the largest observed gains, complemented by independent query sampling.
